\chapter{حديث عبر الهواء}
% full version: chapters/05_serocomputer.tex

\pencilfig{5}{\pencilart{5}}{إلى اليسار الهاتف: سلك من واحد إلى واحد. وإلى اليمين الراديو: رسالة واحدة لكل من حولها دفعة واحدة.}

حتى الآن تحدثت العصبونات بعضها إلى بعض بالهاتف: سلك إلى سلك، بعنوان محدد. والآن نشغّل أول محطة إذاعية: عصبونات \nt{SERO} التي تقذف السيروتونين. ليس منها في الدماغ كله سوى بضع مئات الآلاف، لكن أسلاكها تتشعب في مناطق شاسعة منه.

\section*{أربعة فوارق}

عصبون \nt{SERO} لا يشبه \nt{GLUT}. أولًا، المستقبِل هو من يختار أتكون إشارته ”زائدًا“ أم ”ناقصًا“: للسيروتونين ”أقفال“ من النوعين كليهما. ثانيًا، يعمل من تلقاء نفسه: في اليقظة ينقر بانتظام بضع مرات في الثانية إذا تلقى تغذية خلفية ثابتة، فهو إذن ميقاتية (مترونوم). ثالثًا، هو بطيء: إشارته تفعل فعلها لا في أجزاء من ألف من الثانية، بل في مئات منها. رابعًا، كثيرًا ما يقذف مادته لا في فجوة ضيقة، بل في الحيز المحيط، حيث تنتشر كقطرة حبر في الماء. يحس المستقبِل التركيز، ولا يعرف ممن جاء الجزيء.

\section*{منطق غير موجود}

ميقاتية ذات ”قفل“ كابح هي ”لا“ جاهزة: ما دام السيروتونين غائبًا، يعمل العصبون؛ وإذا ظهر، صمت. عصبون واحد بدل دائرة كاملة. يبدو أن شبكة السيروتونين أقوى منطقيًا من شبكة \nt{GLUT}.

لكن ثمة عقبة. المادة المقذوفة في الحيز تختلط. تبقى الإشارتان متمايزتين فقط إذا أُطلقتا في موضعين يتباعدان أكثر مما تلحق المادة أن تنتشر، أي عشرات الميكرومترات. غير أن لكل عصبون سيروتونين عددًا هائلًا من مواضع الإطلاق، مبعثرة في مناطق واسعة، و”أسلاك“ العصبونات المختلفة تتداخل. لا يمكن تمرير بتات منفصلة عبر شبكة كهذه. لا تستطيع أن تنقل إلا بضع كميات عامة، كراديو يستمع فيه الجميع إلى المحطة نفسها.

\pencilfig{123}{%
% two release sites: far apart the clouds stay separate (two signals), close together they merge (one level)
\foreach \x/\y in {0.9/1.55, 4.1/1.55, 1.9/0, 2.9/0}{
  \foreach \r/\o in {0.95/0.07, 0.7/0.09, 0.45/0.12}{\fill[pcViolet, opacity=\o] (\x,\y) circle (\r);}
  \draw[dotpen=pcViolet, fill=pcViolet] (\x,\y) circle (0.07);}
\draw[arr] (1.95,1.55) -- (3.05,1.55); \draw[arr] (3.05,1.55) -- (1.95,1.55);
\node[lbl, anchor=south, align=center] at (2.5,1.6) {أبعد من\\مدى الانتشار};
\node[lbl, anchor=west] at (5.25,1.55) {إشارتان};
\node[lbl, anchor=west] at (5.25,0) {سحابة واحدة مشتركة};
}{يُقذف السيروتونين في الحيز وينتشر سحابةً. يعطي المصدران إشارتين مختلفتين فقط إذا تباعدا أكثر مما تلحق المادة أن تنتشر. وإذا تقاربا، اندمجت السحابتان في مستوى واحد مشترك.}

\section*{ما الذي تفعله حقًا}

لكمية عامة واحدة يملك هذا النظام كل ما يلزم. تركيز السيروتونين يمهّد النقرات المنفردة إلى مستوى سلس، كما يحوّل أبسط محوّل تدفق النبضات إلى جهد. يدوم هذا المستوى ثواني ثم يذوب وحده: ليس بتًا يجب محوه، بل أثرًا يتلاشى ببطء.

ولعصبونات السيروتونين أيضًا ”أقفال“ على نفسها: السيروتونين المقذوف يكبح مصدره قليلًا. سيرى المهندس في هذا مضخمًا بتغذية راجعة سالبة، ترموستاتًا يحفظ المستوى. لكن التجارب تبين أن هذا الكبح الذاتي موجّه إلى عناوين محددة ولا يسبب إلا توقفات قصيرة، فدور الترموستات لا يزال فرضية أكثر منه حقيقة.

\section*{كم يكلف الانتظار}

أي كمية عامة يحسن أن تبقى بطيئة وواحدة للشبكة كلها؟ مرشح جيد: \emph{الصبر}. تخيّل الاختيار: مكافأة صغيرة الآن، أو كبيرة بعد بضع ثوانٍ. كم أنت مستعد للانتظار يتوقف على سرعة ”انخفاض قيمة“ المستقبل عندك. وفق إحدى الفرضيات، السيروتونين هو من يدير هذا المقبض. وثمة تجارب في صالحها: إذا شُغّلت عصبونات السيروتونين اصطناعيًا، انتظر الفأر المكافأة المؤجلة مدة أطول قبل أن يستسلم.

يقدّم هذا الفصل ثلاث أدوات ستستخدمها كل محطات الدماغ الإذاعية: العصبون الميقاتية، و”السلك“ الذي هو في الحقيقة حجم، والمستوى البطيء بدل النقرات المنفردة.

\fullversion{5}
